\section{Лекция 3. Введение в комбинаторику}

\subsection{Основные понятия комбинаторики}

\begin{definition}
    \textdef{Комбинаторика} — это раздел дискретной математики, 
    изучающий дискретные объекты, множества (комбинаторные конфигурации) 
    и отношения на них. 
\end{definition}

В частности, комбинаторика решает задачи подсчёта 
числа элементов конечных множеств, удовлетворяющих заданным правилам.

\begin{definition}
    \textdef{Комбинаторные конфигурации} — это структурированные 
    определенным образом наборы объектов, составленные из элементов 
    конечных множеств по заданным правилам.
\end{definition}

\begin{example}
    Примерами конфигураций являются перестановки, сочетания,
    размещения и разбиения.
\end{example}

В комбинаторике мы работаем с наборами объектов. Любой \textdef{набор} 
можно классифицировать по двум основным критериям:
\begin{itemize}
    \item \textbf{Упорядоченность:}
    \begin{itemize}
        \item \emph{Неупорядоченный набор}: порядок элементов 
        не имеет значения.
        \item \emph{Упорядоченный набор}: порядок элементов 
        строго зафиксирован (кортеж).
    \end{itemize}
    \item \textbf{Наличие повторений:}
    \begin{itemize}
        \item \emph{Без повторений}: все элементы набора уникальны. 
        Неупорядоченный набор без повторений образует классическое 
        \textdef{множество}.
        \item \emph{С повторениями}: элементы могут встречаться 
        несколько раз. Неупорядоченный набор с повторениями называется 
        \textdef{мультимножеством}.
    \end{itemize}
\end{itemize}

\begin{designation}
    Часто используются сокращения для указания мощности 
    множеств и наборов:
    \begin{itemize}
        \item \textdef{\(n\)-набор} — набор, состоящий из \(n\) элементов;
        \item \textdef{\(n\)-множество} — множество из \(n\) элементов;
        \item \textdef{\(n\)-мультимножество} — мультимножество из 
        \(n\) элементов;
        \item \textdef{\(n\)-кортеж} — упорядоченный набор 
        из \(n\) элементов.
    \end{itemize}
\end{designation}

\begin{definition}
    \textdef{\((n, k)\)-выборка} — это любой \(k\)-набор, составленный из 
    элементов заданного \(n\)-множества.
\end{definition}

\subsection{Базовые принципы}
В данном разделе мы рассматриваем только конечные множества.

\begin{theorem}[\textbf{0. Принцип биективности}]
    Если между двумя множествами существует биективное отображение, 
    то они имеют равное количество элементов.
\end{theorem}

\begin{theorem}[\textbf{1. Принцип сложения}]
    Если множества \(X_1\) и \(X_2\) не пересекаются 
    (\(X_1 \cap X_2 = \varnothing\)), то мощность их объединения равна 
    сумме их мощностей:
    \[
        |X_1 \cup X_2| = |X_1| + |X_2|
    \]
\end{theorem}

\begin{theorem}[\textbf{2. Принцип умножения}]
    Мощность декартова произведения конечных множеств равна произведению 
    их мощностей:
    \[
        |X_1 \times X_2 \times \dots \times X_m| 
        = |X_1| \cdot |X_2| \cdot \dots \cdot |X_m|
    \]
\end{theorem}

\begin{example}
    Пусть существуют 3 дороги из города A в город B, и 2 дороги из города 
    B в город C. Множество маршрутов из A в B обозначим как \(X_1\), 
    из B в C — \(X_2\). Составной маршрут — это элемент декартова 
    произведения \(X_1 \times X_2\). По принципу умножения, общее 
    количество путей из A в C равно \(3 \cdot 2 = 6\).
\end{example}

\begin{theorem}[\textbf{3. Принцип Дирихле}]
    Пусть \(m\) предметов («кроликов») распределены по \(n\) классам 
    («ящикам»).
    \begin{itemize}
        \item Если \(m > n\), то хотя бы в одном ящике окажется строго 
        больше одного предмета (\(> 1\)).
        \item \emph{Обобщенный принцип:} Если \(m > n \cdot k\) (или, в 
        более общем виде, \(m > n \cdot k + c\), где \(c \ge 0\)), то 
        хотя бы в одном ящике окажется строго больше \(k\) предметов.
    \end{itemize}
\end{theorem}

\begin{proof}[Обоснование обобщенного принципа]
    \begin{quote}
        Докажем от противного. Пусть мы распределяем \(m > n \cdot k\) 
        предметов по \(n\) ящикам, и предположим, что в каждом ящике 
        оказалось не более \(k\) предметов. Тогда общее число предметов во 
        всех ящиках не превышает \(n \cdot k\). Однако по условию общее 
        число предметов \(m > n \cdot k\). Мы получили противоречие. 
        Значит, хотя бы в одном ящике находится как минимум \(k + 1\) 
        предмет (то есть строго больше \(k\)).
    \end{quote}
\end{proof}

\begin{example}
    Бросили 5 точек в квадрат размером \(1 \times 1\). Докажем, что 
    найдутся хотя бы 2 точки, расстояние между которыми не превышает 
    \(1/\sqrt{2}\).
    
    \emph{Доказательство:} Разобьем исходный квадрат на 4 равных квадрата 
    размером \(1/2 \times 1/2\). Мы имеем \(m = 5\) точек (кроликов) и 
    \(n = 4\) квадрата (ящика). Так как \(5 > 4\), по принципу Дирихле 
    хотя бы в один из маленьких квадратов попадут не менее 2 точек. 
    Максимальное расстояние между любыми двумя точками внутри квадрата со 
    стороной \(1/2\) достигается на концах его диагонали и равно 
    \(\sqrt{(1/2)^2 + (1/2)^2} = \sqrt{1/4 + 1/4} = \sqrt{1/2} 
    = 1/\sqrt{2}\). Следовательно, расстояние между этими двумя точками не 
    превосходит \(1/\sqrt{2}\).
\end{example}

\begin{figure}[H]
    \centering
    \begin{tikzpicture}[scale=2.5]
        % Внешний квадрат
        \draw[thick] (0,0) rectangle (1,1);
        
        % Разбиение на 4 малых квадрата
        \draw[dashed] (0.5, 0) -- (0.5, 1);
        \draw[dashed] (0, 0.5) -- (1, 0.5);
        
        % Подписи сторон
        \node[left] at (0, 0.5) {\(1\)};
        \node[below] at (0.5, 0) {\(1\)};
        
        % 5 точек (по принципу Дирихле)
        % Квадрат 1 (верх-право): 1 точка
        \fill (0.7, 0.8) circle (1pt);
        % Квадрат 2 (верх-лево): 2 точки (выделены красным)
        \fill[red] (0.2, 0.6) circle (1.2pt);
        \fill[red] (0.4, 0.9) circle (1.2pt);
        % Квадрат 3 (низ-лево): 1 точка
        \fill (0.3, 0.2) circle (1pt);
        % Квадрат 4 (низ-право): 1 точка
        \fill (0.8, 0.3) circle (1pt);
        
        % Линия расстояния между красными точками
        \draw[<->, red, densely dotted, thick] (0.22, 0.62) 
            -- (0.38, 0.88) 
            node[midway, right=2pt, scale=0.9] {\(\le \frac{1}{\sqrt{2}}\)};
            
        % Поясняющий текст
        \node[anchor=west, text width=4cm, align=left] at (1.2, 0.5) {
            \textbf{5} точек (кроликов)\\
            \textbf{4} квадрата (ящика)\\[0.5em]
            \(\Rightarrow\) хотя бы в одном \\
            квадрате \(\ge 2\) точек
        };
    \end{tikzpicture}
\end{figure}

\begin{theorem}[\textbf{4. Принцип включений и исключений}]
    Для двух множеств:
    \[
        |X_1 \cup X_2| = |X_1| + |X_2| - |X_1 \cap X_2|
    \]
    Для трех множеств:
    \begin{align*}
        |X_1 \cup X_2 \cup X_3| = {} & |X_1| + |X_2| + |X_3| - 
        |X_1 \cap X_2| - |X_1 \cap X_3| - |X_2 \cap X_3| \\
        & + |X_1 \cap X_2 \cap X_3|
    \end{align*}
    В общем случае для \(n\) множеств формула имеет вид:
    \[
        |X_1 \cup X_2 \cup \dots \cup X_n| = \sum_{k=1}^{n} (-1)^{k-1} 
        \left( \sum_{1 \le i_1 < i_2 < \dots < i_k \le n} 
        |X_{i_1} \cap \dots \cap X_{i_k}| \right)
    \]
\end{theorem}

\begin{proof}[Доказательство принципа включений и исключений]
    \begin{quote}
        Покажем, что каждый элемент множества, находящийся в объединении 
        \(X_1 \cup \dots \cup X_n\), будет подсчитан правой частью 
        формулы ровно 1 раз. 
        
        Рассмотрим произвольный элемент \(x\), который принадлежит ровно 
        \(m\) множествам из данных \(n\) (где \(1 \le m \le n\)). 
        В левой части формулы (в объединении) этот элемент 
        считается ровно 1 раз. 
        
        Посмотрим, какой вклад этот элемент дает в правую часть:
        \begin{itemize}
            \item При \(k=1\) (суммирование одиночных множеств) он будет 
            посчитан \(\binom{m}{1}\) раз.
            \item При \(k=2\) (суммирование попарных пересечений) он будет 
            посчитан \(\binom{m}{2}\) раз, причем со знаком минус.
            \item В общем случае, в сумме пересечений по \(k\) множествам он 
            встретится \(\binom{m}{k}\) раз со знаком \((-1)^{k-1}\).
        \end{itemize}
        Следовательно, общий вклад элемента \(x\) в правую часть равен:
        \[
            S = \binom{m}{1} - \binom{m}{2} + \binom{m}{3} - \dots + 
            (-1)^{m-1} \binom{m}{m}
        \]
        Воспользуемся биномом Ньютона для разложения \((1 - 1)^m\):
        \[
            0 = (1 - 1)^m = \binom{m}{0} - \binom{m}{1} + \binom{m}{2} 
            - \dots + (-1)^m \binom{m}{m}
        \]
        Так как \(\binom{m}{0} = 1\), перенесем все остальные члены в 
        левую часть:
        \[
            1 = \binom{m}{1} - \binom{m}{2} + \dots + (-1)^{m-1} \binom{m}{m}
        \]
        То есть \(S = 1\). Значит, каждый элемент действительно 
        учитывается формулой ровно один раз.
    \end{quote}
\end{proof}

\subsubsection{Частные случаи применения принципов}

\textbf{Частный случай №1: Слова в алфавите.} Пусть \(\Sigma\) — некоторый 
алфавит. Множество всех возможных слов в алфавите \(\Sigma\) длины \(k\) 
обозначается \(\Sigma^k\). По принципу умножения:
\[
    |\Sigma^k| = |\Sigma|^k
\]
В частности, для бинарного алфавита \(\Sigma = \{0, 1\}\) количество 
слов длины \(k\) равно \(2^k\).

\vspace{2 em}

\textbf{Частный случай №2: Подмножества конечного множества.} Пусть 
\(X = \{x_1, \dots, x_n\}\) — конечное множество. Мощность его булеана 
(множества всех подмножеств) \(\mathcal{P}(X)\) равна \(2^{|X|}\). 

\vspace{1 em}

Покажем это с помощью бинарных векторов длины \(n\) (\(\B^n\)): Каждому 
подмножеству можно сопоставить вектор, где на \(i\)-й позиции стоит \(1\), 
если элемент \(x_i\) входит в подмножество, и \(0\) — если не входит. 

Например, если \(n=5\), то подмножеству \(\{x_1, x_4, x_5\}\) соответствует 
бинарный вектор \(10011\) (в порядке позиций \(x_1, x_2, x_3, x_4, x_5\)). 
Векторов длины \(n\) ровно \(2^n\), значит и подмножеств \(2^n\).

\begin{figure}[H]
    \centering
    \begin{tikzpicture}[scale=0.65]
        % Подписи элементов множества сверху
        \node[above] at (0.5, 1.1) {\(x_1\)};
        \node[above] at (1.5, 1.1) {\(x_2\)};
        \node[above] at (2.5, 1.1) {\(x_3\)};
        \node[above] at (3.5, 1.1) {\(x_4\)};
        \node[above] at (4.5, 1.1) {\(x_5\)};
        
        % Отрисовка ячеек бинарного вектора
        \draw[thick] (0,0) rectangle (1,1) node[midway] {\(1\)};
        \draw[thick] (1,0) rectangle (2,1) node[midway] {\(0\)};
        \draw[thick] (2,0) rectangle (3,1) node[midway] {\(0\)};
        \draw[thick] (3,0) rectangle (4,1) node[midway] {\(1\)};
        \draw[thick] (4,0) rectangle (5,1) node[midway] {\(1\)};
        
        % Поясняющий текст слева и справа
        \node[anchor=east] at (-0.5, 0.5) {Подмножество 
            \(\{x_1, x_4, x_5\} \implies\)};
        \node[anchor=west] at (5.5, 0.5) {\(\in \B^5\)};
    \end{tikzpicture}
\end{figure}

\subsection{Перестановки, размещения и сочетания}

\subsubsection{Перестановки}
\begin{definition}
    \textdef{Перестановка} конечного множества \(X = \{x_1, \dots, x_n\}\) 
    — это биективное отображение из множества \(X\) в себя (\(X \to X\)).
\end{definition}

\begin{designation}
    Часто перестановку записывают в виде матрицы из двух строк:
    \[
        \begin{pmatrix} 
            x_1 & x_2 & \dots & x_n \\ 
            x_{i_1} & x_{i_2} & \dots & x_{i_n} 
        \end{pmatrix}
    \]
    Или, при фиксированном порядке элементов, кратко кортежем 
    \((x_{i_1}, x_{i_2}, \dots, x_{i_n})\).
\end{designation}

\textbf{Частный случай №3:} Число всех возможных перестановок 
\(n\)-множества равно факториалу: \(n!\).

\begin{theorem}[\textbf{Формула Стирлинга}]
    Для больших значений \(n\) факториал можно аппроксимировать следующей 
    формулой:
    \[
        n! \sim \sqrt{2 \pi n} \left(\frac{n}{e}\right)^n
    \]
\end{theorem}

\begin{proof}[Эскиз доказательства формулы Стирлинга]
    \begin{quote}
        Рассмотрим натуральный логарифм факториала:
        \[
            \ln(n!) = \ln(1 \cdot 2 \cdot \dots \cdot n) = \sum_{k=1}^{n} \ln k
        \]
        Эту сумму можно приблизить определенным интегралом от функции \(\ln x\):
        \[
            \sum_{k=1}^{n} \ln k \approx \int_{1}^{n} \ln x \, dx
        \]
        Первообразная функции \(\ln x\) равна \(x \ln x - x\). Вычисляя 
        интеграл по формуле Ньютона-Лейбница, получаем:
        \[
            \int_{1}^{n} \ln x \, dx = (n \ln n - n) - (1 \ln 1 - 1) = 
            n \ln n - n + 1
        \]
        Отбрасывая младшие члены, получаем \(\ln(n!) \approx n \ln n - n\). 
        Потенцируя обе части, находим главную часть асимптотики:
        \[
            n! \approx e^{n \ln n - n} = (e^{\ln n})^n \cdot e^{-n} = 
            n^n \cdot e^{-n} = \left(\frac{n}{e}\right)^n
        \]
        Более точный математический анализ (с использованием формулы Тейлора 
        для оценки ошибки интегрирования и формулы Валлиса для нахождения 
        константы) дает дополнительный множитель \(\sqrt{2 \pi n}\), что 
        окончательно приводит к формуле Стирлинга.
    \end{quote}
\end{proof}


\subsubsection{Размещения (Arrangements)}
\begin{definition}
    \textdef{Размещением (без повторений)} \(k\) элементов из 
    \(n\)-множества называется упорядоченная \((n, k)\)-выборка 
    без повторений.
\end{definition}

\begin{designation}
    Обозначается \(A_n^k\) или \((n)_k\).
\end{designation}
Количество таких размещений вычисляется по формуле:
\[
    A_n^k = n \cdot (n-1) \cdot (n-2) \cdot \dots \cdot (n-k+1) = 
    \frac{n!}{(n-k)!}
\]

\begin{definition}
    \textdef{Размещением (с повторениями)} \(k\) элементов из 
    \(n\)-множества называется упорядоченная \((n, k)\)-выборка 
    с возможными повторениями элементов.
\end{definition}

\begin{designation}
    Обозначается \(\bar{A}_n^k\).
\end{designation}
На каждую из \(k\) позиций мы можем выбрать любой из \(n\) элементов, 
поэтому:
\[
    \bar{A}_n^k = n^k
\]


\subsubsection{Сочетания (Combinations)}
\begin{definition}
    \textdef{Сочетанием (без повторений)} \(k\) элементов из 
    \(n\)-множества называется неупорядоченная \((n, k)\)-выборка 
    без повторений.
\end{definition}

\begin{designation}
    Обозначается \(C_n^k\) или \(\binom{n}{k}\).
\end{designation}

Количество сочетаний без повторений выражается через размещения:
\[
    C_n^k = \frac{A_n^k}{k!}
\]

\begin{remark}[Связь через отношение эквивалентности]
    Сочетания можно рассматривать через факторизацию размещений по 
    отношению эквивалентности. 
    Пусть \(X\) — множество всех размещений без повторений (\(k\)-кортежей) 
    из \(n\)-множества. Мощность этого множества равна \(A_n^k\).
    Введем отношение эквивалентности \(\sim\): два размещения считаются 
    эквивалентными, если они состоят из одного и того же набора элементов 
    (отличаются лишь их порядком). 
    Так как переставить \(k\) элементов можно ровно \(k!\) способами, 
    каждый класс эквивалентности содержит ровно \(k!\) элементов. 
    Поскольку классы не пересекаются, общее количество таких классов равно 
    \(A_n^k / k!\). Сам класс эквивалентности представляет собой 
    неупорядоченный набор, то есть \textdef{сочетание}.
\end{remark}


\begin{definition}
    \textdef{Сочетанием (с повторениями)} \(k\) элементов из 
    \(n\)-множества называется неупорядоченная \((n, k)\)-выборка 
    с повторениями.
\end{definition}

\begin{designation}
    Обозначается \(\bar{C}_n^k\).
\end{designation}

Количество таких сочетаний равно:
\[
    \bar{C}_n^k = C_{n+k-1}^{n-1} = C_{n+k-1}^k
\]

\begin{proof}[Обоснование формулы через перегородки 
(метод шаров и перегородок)] ~
    \begin{quote}
        Любую неупорядоченную \((n, k)\)-выборку с повторениями можно 
        представить как задачу размещения \(k\) неразличимых шаров по \(n\) 
        различимым ящикам (каждый ящик соответствует своему типу элемента).
        
        Выстроим шары в ряд. Чтобы разделить их на \(n\) ящиков, нам 
        понадобится поставить между ними \(n-1\) перегородку. 
        Таким образом, мы формируем последовательность из символов 
        двух типов: \(k\) шаров (обозначим их «\(0\)») и \(n-1\) 
        перегородок (обозначим их «\(1\)»). Общая длина такой 
        последовательности (бинарного вектора) составит \(k + n - 1\).
        
        Например, вектор \(0110010\) для \(n=4\) ящиков (3 перегородки) 
        и \(k=4\) шаров означает: в первый ящик попал 1 шар, во второй 
        — 0 шаров, в третий — 2 шара, в четвертый — 1 шар.

        \begin{figure}[H]
            \centering
            \begin{tikzpicture}[scale=0.75] 
                % Отрисовка ячеек вектора
                \draw[thick] (0,0) rectangle (1,1) node[midway] {\(0\)};
                \draw[thick] (1,0) rectangle (2,1) node[midway] {\(1\)};
                \draw[thick] (2,0) rectangle (3,1) node[midway] {\(1\)};
                \draw[thick] (3,0) rectangle (4,1) node[midway] {\(0\)};
                \draw[thick] (4,0) rectangle (5,1) node[midway] {\(0\)};
                \draw[thick] (5,0) rectangle (6,1) node[midway] {\(1\)};
                \draw[thick] (6,0) rectangle (7,1) node[midway] {\(0\)};
                
                % Имитация фигурных скобок снизу (номера ящиков)
                \draw (0.1, -0.2) -- (0.1, -0.4) 
                    -- (0.9, -0.4) -- (0.9, -0.2);
                \node at (0.5, -0.7) {\(1\)};
                
                \draw (3.1, -0.2) -- (3.1, -0.4) 
                    -- (4.9, -0.4) -- (4.9, -0.2);
                \node at (4.0, -0.7) {\(3\)};
                
                \draw (6.1, -0.2) -- (6.1, -0.4) 
                    -- (6.9, -0.4) -- (6.9, -0.2);
                \node at (6.5, -0.7) {\(4\)};
                
                % Подписи сверху (перегородки)
                \draw[<-] (1.5, 1.1) -- (1.5, 1.8);
                \draw[<-] (2.5, 1.1) -- (2.5, 1.8);
                \draw[<-] (5.5, 1.1) -- (5.5, 1.8) 
                    node[above, xshift=-1cm] {перегородки (\(n-1\) шт.)};
                
                % Дополнительный текст справа
                \node[anchor=west, text width=5cm] at (7.5, 0.5) {
                    \(\B^{n+k-1}\) при условии,\\
                    что кол-во единиц \(= n-1\)
                };
            \end{tikzpicture}
        \end{figure}
        
        Чтобы найти количество всех таких последовательностей (а значит, и 
        число сочетаний), достаточно из общего числа позиций \((n+k-1)\) 
        выбрать \((n-1)\) позиций для расстановки единиц (перегородок). 
        Это делается способами:
        \[
            C_{n+k-1}^{n-1} \quad \text{или, что эквивалентно,} \quad 
            C_{n+k-1}^k
        \]
        Таким образом, задача сводится к выбору бинарных векторов из 
        \(\B^{n+k-1}\) при условии, что количество единиц равно \(n-1\).
    \end{quote}
\end{proof}

\subsubsection{Классификация (n, k)-выборок}

Подводя итог, приведем сводку формул для вычисления \((n, k)\)-выборок:

\begin{figure}[H]
    \centering
    \begin{tikzpicture}[>=Stealth, thick, xscale=1.0, yscale=1.2]
        % Корневой узел
        \node (root) at (0, 0) {\((n, k)\)-выборка};
        
        % Уровень 1
        \node[align=center] (upor) at (3.5, 1.5) 
            {упорядоченная\\(размещение)};
        \node[align=center] (neupor) at (3.5, -1.5) 
            {неупорядоченная\\(сочетание)};
        
        \draw[->] (root.east) -- (upor.west);
        \draw[->] (root.east) -- (neupor.west);
        
        % Уровень 2
        \node[anchor=west] (ub) at (6.5, 2.2) {без повторений: 
            \(A_n^k = \frac{n!}{(n-k)!}\)};
        \node[anchor=west] (us) at (6.5, 0.8) {с повторениями: 
            \(\bar{A}_n^k = n^k\)};
        
        \node[anchor=west] (nb) at (6.5, -0.8) {без повторений: 
            \(C_n^k = \frac{n!}{k!(n-k)!}\)};
        \node[anchor=west] (ns) at (6.5, -2.2) {с повторениями: 
            \(\bar{C}_n^k = \frac{(n+k-1)!}{k!(n-1)!}\)};
        
        \draw[->] (upor.east) -- (ub.west);
        \draw[->] (upor.east) -- (us.west);
        \draw[->] (neupor.east) -- (nb.west);
        \draw[->] (neupor.east) -- (ns.west);
        
        % Дополнительные обозначения справа
        \node[anchor=west] at (11.5, 2.2) {\(A_n^k = (n)_k\)};
        \node[anchor=west] at (11.5, -0.8) {\(C_n^k = \binom{n}{k}\)};
    \end{tikzpicture}
\end{figure}

\subsection{Пример решения задачи}

\begin{problem}
    Дано число 720. Сколько всего целых положительных делителей 
    у этого числа?
\end{problem}

\begin{proof}[Решение] ~
    \begin{quote}
        Разложим число 720 на простые множители:
        \[
            720 = 2^4 \cdot 3^2 \cdot 5^1
        \]
        Любой делитель этого числа будет иметь вид 
        \(2^a \cdot 3^b \cdot 5^c\).
        Найдем, сколько вариантов существует для выбора каждой из степеней:
        \begin{itemize}
            \item \(a \in \{0, 1, 2, 3, 4\}\) — 5 вариантов;
            \item \(b \in \{0, 1, 2\}\) — 3 варианта;
            \item \(c \in \{0, 1\}\) — 2 варианта.
        \end{itemize}
        По правилу умножения общее количество делителей составит:
        \[
            5 \cdot 3 \cdot 2 = 30 \text{ вариантов}
        \]
    \end{quote}
\end{proof}